% !TEX root = ../manuscript.tex

\section{End-to-end uncertainty quantification workflow}

The framework connects uncertainties that are commonly treated in isolation:
fault--seal geology, unresolved fault-core architecture, spatially variable
full-fault properties, and field-scale \coTwo{} migration and containment
(Fig.~\ref{fig:end-to-end-workflow}). It has three linked components---local
fault-property modeling, full-fault property sampling, and field-scale storage
simulation. This hierarchy allows broad geologic effects to be distinguished
from residual stochastic variability conditional on a prescribed geology.

At the local scale, the 162 geologic scenarios specify structural and
stratigraphic conditions within six throw windows. For every geology--window
combination, PREDICT generates 2,000 stochastic three-dimensional fault-core
architectures and their joint anisotropic permeability. Exact realization
replay and physics-based upscaling then provide internally consistent
permeability, porosity, capillary-pressure, and relative-permeability
properties.

At the full-fault scale, selected local realizations populate six windows
along dip and 87 slices along strike. Phase 1 uses 12 conditionally independent
full-fault realizations per geology to characterize stochastic variability,
while a full-distribution medoid and coherent low- and high-state fields serve
as separate deterministic benchmarks. Reproducible manifests preserve the
selected realization, seed, code version, and property provenance for every
fault element.

At the field scale, JutulDarcy propagates these heterogeneous fault properties
through a roughly two-million-cell offshore Texas storage model over 1,000
years. The balanced ensemble identifies geologic controls and within-geology
variability in migration and containment, and targeted enrichment resolves
conditional outcome distributions for geologies that exhibit mixed behavior.
The study design and assumptions needed to interpret the ensemble are
summarized in Materials and Methods. Complete sampling, replay, upscaling,
simulation, validation, and analysis procedures are provided in the SI
Appendix.
